\documentclass[12pt,a4paper]{article}

\usepackage[a4paper,margin=1in]{geometry}
\usepackage[T1]{fontenc}
\usepackage[utf8]{inputenc}
\usepackage{newtxtext,newtxmath}
\usepackage[hidelinks]{hyperref}
\usepackage{setspace}
\usepackage[ngerman]{babel}
\usepackage{changepage}

\setlength{\parskip}{0.5\baselineskip}
\setlength{\parindent}{0pt}

\pagestyle{empty}
\onehalfspacing

\begin{document}

\textbf{Quin Barnard}\\
\textbf{27902277}\\

Der \underline{volkstümliche} Humor in Südafrika ist dem deutschen ähnlich.\\
Jener wird oft von Ausländer für eine unantsändige \underline{Gepräge} gehalten.\\
Meiner Meinung nach zeigt er auf nichts als unseren unsterblichen \underline{Drang} nach Scharfsinn.\\
Trotzdem, wir müssen vielleicht lernen nicht jede \underline{Hemmung} zu zersprengen.\\
Diese Sätze \underline{widme} ich meiner Mutter.\par

Die \underline{Beredsamkeit} strebt jeder Schriftsteller auf eigene Weise nach.\\
Am besten wird eine Gleichgewicht zwischen dieser und \underline{Prägnanz} gefunden.\\
Beethoven ofenbarte beide Tugenden als er sagte er volle dem \underline{Schicksal} in den Rachen greifen.\\
Taubheit erlaubte er nichs als ein kleiner \underline{Stachel} zu sein.\\
Der Wille als der \underline{unmittelbar} menschliche Eigenschaft wird hier gesehen.\par

\underline{Märchenhaft} empfinde ich die Bibliothek der Stellenbosch Universität.\\
In jenem Ort bilden sich die \underline{Gebildeten}; lehren sich die Lehrern!\\
Die Heiterkeit der Bücher löst bei mir jede \underline{Schwermut} auf.\\
Nichts ist der \underline{Reize} des Wissens ähnlich.\\
Obwohl sie mir fast \underline{unwiderstehlich} hetzt, kann ich leider nicht viel Zeit dort verbringen.\par

\end{document}
